% Part: first-order-logic
% Chapter: introduction
% Section: formulas

\documentclass[../../../include/open-logic-section]{subfiles}

\begin{document}

\olfileid[pt]{fol}{int}{fml}

\section{\usetoken{P}{formula}}

Eis a abordagem que usaremos para especificar rigorosamente as
!!{sentence}s da lógica de primeira ordem e para lidar com as questões
decorrentes do uso de !!{variable}s. Primeiro definimos um conjunto
\emph{diferente} de expressões: !!{formula}s. Uma vez feito isso,
podemos considerar o papel que as !!{variable}s desempenham nelas---e,
com base em algumas outras ideias, nomeadamente as de !!{variable}s
``livres'' e ``ligadas'', podemos definir o que é uma !!{sentence}
(nomeadamente, uma !!{formula} sem !!{variable}s livres). Fazemos
isso não apenas porque torna a definição de ``!!{sentence}'' mais
gerenciável, mas também porque será crucial para a maneira como
definimos a noção semântica de satisfação.

Vamos definir ``!!{formula}'' para uma linguagem simples de primeira
ordem, contendo apenas um único !!{predicate}~$\Obj P$ e um único
!!{constant}~$\Obj a$, e apenas os símbolos lógicos $\lnot$, $\land$
e~$\lexists$. Nossas definições completas serão muito mais gerais:
permitiremos infinitos !!{predicate}s e !!{constant}s. De
fato, também consideraremos !!{function}s que podem ser combinados com
!!{constant}s e !!{variable}s para formar ``termos''. Por enquanto,
$\Obj a$ e as variáveis serão nossos únicos termos.  Precisamos de
infinitas !!{variable}s.  Oficialmente usaremos os símbolos
$\Obj v_0$, $\Obj v_1$, \dots, como variáveis.

\begin{defn}
O conjunto de \emph{!!{formula}s}~$\Frm$ é definido da seguinte forma:
\begin{enumerate}
\item\ollabel{fmls-atom} $\Atom{\Obj P}{\Obj a}$ e $\Atom{\Obj
  P}{\Obj v_i}$ são !!{formula}s ($i \in \Nat$).

\tagitem{prvNot}{\ollabel{fmls-not}Se $!A$ é uma !!{formula}, então $\lnot !A$ é
  uma !!{formula}.}{}

\tagitem{prvAnd}{Se $!A$ e $!B$ são !!{formula}s, então $(!A \land
  !B)$ é uma !!{formula}.}{}

\tagitem{prvEx}{\ollabel{fmls-ex}Se $!A$ é uma !!{formula} e $x$ é uma !!{variable},
  então $\lexists[x][!A]$ é uma !!{formula}.}{}

\tagitem{limitClause}{\ollabel{fmls-limit}Nada mais é uma !!{formula}.}{}
\end{enumerate}
\end{defn}

\olref{fmls-atom} nos diz que $\Atom{\Obj P}{\Obj a}$ e
$\Atom{\Obj P}{\Obj v_i}$ são !!{formula}s, para qualquer $i \in
\Nat$. Estas são as chamadas !!{formula}s \emph{atômicas}. Elas nos
dão algo com que começar.  As outras cláusulas nos dão maneiras de
formar novas !!{formula}s a partir daquelas que já formamos. Assim, por
exemplo, por \olref{fmls-not}, obtemos que $\lnot \Atom{\Obj P}{\Obj
v_2}$ é uma !!{formula}, já que $\Atom{\Obj P}{\Obj v_2}$ já é uma
!!{formula} por \olref{fmls-atom}. Então, por \olref{fmls-ex},
obtemos que $\lexists[\Obj v_2][\lnot \Atom{\Obj P}{\Obj v_2}]$ é
outra !!{formula}, e assim por diante.  \olref{fmls-limit} nos diz
que \emph{somente} cadeias que podemos formar desta maneira contam
como !!{formula}s. Em particular, $\lexists[\Obj v_0][\Atom{\Obj
P}{\Obj a}]$ e $\lexists[\Obj v_0][\lexists[\Obj v_0][\Atom{\Obj
P}{\Obj a}]]$ \emph{contam} como !!{formula}s, e $(\lnot \Atom{\Obj
P}{\Obj a})$ não conta, por causa dos parênteses externos
supérfluos.

Esta maneira de definir !!{formula}s é chamada de \emph{definição
indutiva}, e nos permite provar coisas sobre !!{formula}s usando uma
versão de prova por indução chamada \emph{indução estrutural}. Estas
são discutidas de maneira geral em \olref[mth][ind][idf]{sec} e
\olref[mth][ind][sti]{sec}, que você deve revisar antes de mergulhar
nas provas mais adiante. Basicamente, a ideia é que, se você quer dar
uma prova de que algo é verdadeiro para todas as !!{formula}s, mostra
primeiro que é verdadeiro para as !!{formula}s atômicas e então que
\emph{se} for verdadeiro para qualquer !!{formula}~$!A$ (e~$!B$),
\emph{também} é verdadeiro para $\lnot !A$, $(!A \land !B)$ e
$\lexists[x][!A]$. Por exemplo, isso prova que é verdadeiro para
$\lexists[\Obj v_2][\lnot \Atom{\Obj P}{\Obj v_2}]$: da primeira parte
você sabe que é verdadeiro para a !!{formula} atômica~$\Atom{\Obj
P}{\Obj v_2}$. Então obtém que é verdadeiro para $\lnot \Atom{\Obj
P}{\Obj v_2}$ pela segunda parte, e novamente que é verdadeiro para
$\lexists[\Obj v_2][\lnot \Atom{\Obj P}{\Obj v_2}]$ em si. Como todas
as !!{formula}s são geradas indutivamente a partir de !!{formula}s
atômicas, isso funciona para qualquer uma delas.

\end{document}
